%\section{\label{sec:related}Related Work}

\section{\label{sec:nas}Related Work}
%The goal of NAS is to find a model architecture $\m$ that achieves a high reward $R(\m)$ on a particular task. Recent approaches such as~\citet{nas,nas_module,nas_dqn,bello2017neural,nas_module_dqn} used reinforcement learning to search the combinatorial space of architectures, attempting to model a parameterized policy $\pi(\m; \theta)$ of $\m$ to maximize the expected reward $J(\theta) = \mathbb{E}_{\m \sim \pi(\m; \theta)} \left[ R(\m) \right]$.

%There are two main approaches for designing the search space. The first approach, which we call \emph{macro search space design}, designs the whole architecture all at once. For instance, if we were to design a convolutional network, we would choose how many filters to use at each layer, the filter sizes, how to set up the skip connections, etc. Macro design is used by~\cite{nas,nas_dqn,deep_arc,smash_net,conv_fabrics}. The second approach, which we call \emph{micro search space design}, aims to search for smaller modules which are then replicated and assembled to form the final model. Micro design is used by~\cite{nas_module,nas_module_dqn}. We note that~\cite{nas} also use micro design when designing the RNN cell.

%An architecture $\m$ is often represented as a sequence, with a grammar and meaning specified by the NAS algorithm.  Because modeling variable length sequences is necessary, the reinforcement learning policy $\pi(\m; \theta)$ is often parametrized with a recurrent neural network. $J(\theta)$ is then learned with various reinforcement learning algorithms such as REINFORCE~\citep{nas,reinforce}, Proximal Policy Optimization~\citep{nas_module,ppo}, and Q-learning~\citep{nas_dqn,nas_module_dqn}.

%NAS methods are typically time consuming and computationally expensive. Most significantly,~\cite{nas} and~\cite{nas_module} respectively train in $3$-$4$ weeks and $3$-$4$ days, using $800$ and $450$ GPUs concurrently at any time during their training. A summary from previous NAS approaches can be found in Table~\ref{tab:cifar10_times} in Appendix~\ref{sec:cifar10_details}. Such computational expense of NAS methods is due to the fact that each architecture $\m$ is trained from scratch to compute $R(\m)$. After training $\m$, only the reward $R(\m; \omega^{\text{opt}}_\m)$ is utilized, while the learned parameters of $\omega^{\text{opt}}_\m$, which could be reused, are thrown away.

There is growing interest in improving the efficiency of neural architecture search. Concurrent to our work are the promising ideas of using learning curve prediction to skip bad models~\citep{Baker17}, predicting the accuracies of models before training~\citep{Deng17}, using iterative search method for architectures of growing complexity~\citep{Liu17,Elsken17}, or using hierarchical representation of architectures~\citep{Liu17b}. Our method is also inspired by the concept of weight inheritance in neuro-evolution, which has been demonstrated to have positive effects at scale~\citep{Real17}.

Closely related to our method are other recent approaches that avoid training each architecture from scratch, such as convolutional neural fabrics -- ConvFabrics~\citep{conv_fabrics} and SMASH~\citep{smash_net}. These methods are more computationally efficient than standard NAS. However, the search space of ConvFabrics is not flexible enough to include novel architectures, e.g. architectures with arbitrary skip connections as in~\cite{nas}. Meanwhile, SMASH can design interesting architectures but requires a hypernetwork~\citep{hyper_net} to generate the weights, conditional on the architectures. While a hypernetwork can efficiently \textit{rank} different architectures, as shown in the paper, the real performance of each network is different from its performance with parameters generated by a hypernetwork. Such discrepancy in SMASH can cause misleading signals for reinforcement learning. Even more closely related to our method is PathNet~\citep{fernando2017pathnet}, which uses evolution to search for a path inside a large model for transfer learning between Atari games.


%These methods are either more expensive or design less accurate models compared to our methods.  

% In this work, we propose a novel approach that combines the strengths of both ConvFabrics and SMASH. Unlike SMASH, we do not use a hypernetwork to generate $\omega_\m$ for each $\m$, but instead keep a shared set of parameters $\omega$ among all architectures $\m$. Unlike ConvFabrics, we design a different mechanism to specify computing paths in our model, leading to a very flexible search space. In fact, our search space is larger than the search space of SMASH as well as the search space described in~\cite{nas}. As a result, our method can efficiently discover novel architectures that achieve better accuracy than SMASH.

% Our work is related to several directions in the existing literature. First, we
% show that how all existing approaches for NAS can be viewed as an instance of
% the Bayesian framework that we describe in Section~\ref{sec:methods}. Then, we
% review existing attempts to establish theoretical grounds for NAS.
% 
% \subsection{\label{sec:existing}Approaches for Neural Architecture Search}
% Existing methods for NAS can be classified into two categories. The first
% category of algorithms attempt to design the the architecture of the whole
% network, e.g.  how many channels to use at each convolutional layers, which
% layers should be sent to which layers via residual connections~\cite{res_net},
% etc. Example algorithms of this approach include~\cite{nas} and
% ~\cite{nas_dqn}.  The second category aims to design several particular modules
% and replicate it multiple times to form a model.  Examples of this category
% include~\cite{nas_module}, where the authors attempt to learn convolutional
% modules and pooling modules to build their networks; and~\cite{nas_module_dqn},
% where the authors use a deep Q-network instead of policy gradient to design
% their convolutional modules.~\cite{nas} also contribute an approach to learn
% the formula for a recurrent neural network cell, which is replicated through
% multiple time steps.
% 
% All the approaches mentioned in the previous paragraphs strives to model the
% distribution above possible architectures in their corresponding search space,
% i.e. $q(\m)$ in Eqn.~\ref{eqn:models_and_params}. The parameter distribution
% $q(\theta_\m | \m)$ is not explicitly modeled, and $\theta_\m$ is learned by
% minimizing a cross-entropy loss objective $\theta_\m = \arg\min_{\theta}
% \mathcal{L(\m, \theta)}$.
% 
% Similar approaches such as~\cite{fractal_net,conv_fabrics} and~\cite{smash_net}
% tries to anneal the running time by designing different choices of $q(\m)$ and
% $q(\theta_\m | \m)$.
% 
% Another inspiring method is to derive architectures via Reproduced Kernel
% Hilbert Space (RKHS)~\cite{kernel_nas}.
% 
% 
% \subsection{\label{sec:theory}Existing Theory}
% There are extremely few theory papers on this NAS.
% Recently,~\cite{non_parametric_net} shows that the optimization process has a
% global minimal. More recently,~\cite{ada_net} prove that the Rademacher
% complexity of the set of models is not significantly larger than that of a
% single model, suggesting an easy optimization problem.


